\begin{papercard}{campbell2013}
\cf{Stages} \stagetag{C} \stagetag{I}. A procedure for building a reliable coding scheme when \emph{one knowledgeable coder} will code everything: standardise the text units first (the principal investigator brackets them), then raise discriminant capability using independent coding (reliability) and \emph{negotiated agreement}. Developed on 101 policy-expert interviews, but explicitly offered for other qualitative data (documents, images, video).

\cf{Problem} Interview-coding reliability is rarely reported or explained. Only 69\,\% of 137 content-analysis articles contained any information on intercoder reliability \pc{campbell2013}{3}{Only 69 per- cent of these articles contained any information on intercoder reliability}. Semistructured interviews raise a distinct \emph{unitization problem}: appropriate blocks of text are not given, and answers are rambling and often need several codes at once \pc{campbell2013}{4}{this unitization prob- lem—that is, identifying appropriate blocks of text for a particular code or codes}.

\cf{Key distinctions} Three types of reliability (Krippendorff): stability, accuracy, reproducibility. Here: \emph{reproducibility}. \emph{Intercoder reliability}: equally capable coders working in isolation pick the same code for the same unit. \emph{Intercoder agreement}: coders reconcile discrepancies through discussion, e.g.\ when one is more knowledgeable \pc{campbell2013}{4}{Intercoder agreement requires that two or more coders are able to reconcile through discussion}. Three-stage process: (1) develop the scheme to the highest reliability on a sample; (2) adjudicate remaining disagreements by negotiation; (3) deploy on the full set by one knowledgeable coder. Two coders are needed for stages 1--2 \pc{campbell2013}{5}{We argue that coding in-depth semistructured interviews should typically involve a three-stage process}.

\cf{Developing the scheme} The PI drafted several dozen grouped, defined codes and a research assistant was trained. Both coded a random transcript, and the scheme was found too cumbersome and unreliable, so codes were dropped or merged, definitions made precise with examples, and the test repeated. A \emph{different} transcript was used each trial. After 11 transcripts (about 10\,\%) they were satisfied \pc{campbell2013}{7}{In every subsequent coding trial, we used a different full-length interview transcript}. Simplification must not destroy nuance. The final scheme used code families, with 60 codes (9 primary, 51 secondary), and each transcript was coded \emph{twice}, once per half of the families, to reduce cognitive load \pc{campbell2013}{8}{It included 60 codes (9 primary and 51 secondary) excluding the names of organizations}, \pc{campbell2013}{15}{Hence, we split the coding scheme into two sets of code families roughly equal in size}.

\cf{Unitization solution} Meaning units (not sentence or page units, since screen sizes vary). The PI codes a transcript in Atlas.ti, saves it, makes a copy, \emph{removes the codes but keeps the brackets}, and gives it to the assistant, who codes the already bracketed units \pc{campbell2013}{11}{Then on a copy he removed all the codes but not the brackets}. Why the PI unitises: it requires subject expertise. Possible bias: unitising may alert coders to codable text, which the authors consider minor because disagreements were mostly about how much text, not which code. A suggested check is to unitise the same transcript narrowly and broadly and compare reliability \pc{campbell2013}{20}{one way to explore this would be to unitize the same transcript in two ways}.

\cf{Results of the two stages} Reliability plateaued at 54\,\% (primary + secondary) and 65\,\% (primary only), explained by exploratory design, long multi-theme units, a complex scheme and an uneven-knowledge coder pair \pc{campbell2013}{12}{finally reached a plateau of 54 percent reliability on average for the primary and secondary codes combined}. After negotiating, \emph{agreement} was 96\,\% and 97\,\%; 91\,\% of disagreements were reconciled, 9\,\% remained irreconcilable, and the assistant deferred to the PI 69\,\% of the time \pc{campbell2013}{13}{After nego- tiating discrepancies, we reached 96 percent intercoder agreement for the primary and secondary codes combined}. Power dynamics may influence deference, so tracking who deferred and discussing it openly is advised \pc{campbell2013}{13}{The potential difficulty with this approach, of course, is the interpersonal dynamics}. Deployment: the PI coded the codes needing expertise and the assistant the straightforward ones.

\cf{Statistic and thresholds} Reliability per code = units where \emph{both} used the code / units where \emph{either} used it, overall = total agreements / (agreements + disagreements) \pc{campbell2013}{15}{we divided the number of times that all coders used it in the same text unit by the number of times that any coder used it}. Krippendorff's $\alpha$ was rejected here because codes were not equiprobable and the coders were not equally qualified. Other reasons: many codes (less chance agreement), multi-coding of units, and no intent to generate statistical variables. No universal threshold exists; aim 80--90\,\% (Miles \& Huberman) and relax for exploratory work, but ``not so low that the findings can no longer be taken seriously'' \pc{campbell2013}{17}{Our initial aim was to achieve a discriminant coding capability of 80 percent to 90 percent}.

\cf{Other practices (seven projects)} Purely inductive single-coder projects ignoring reliability; double-coding everything then reconciling; subset checks; gold-standard comparison against the PI's master coding with 90\,\% needed (coders failing after the fourth iteration were dismissed); first-layer broad codes by assistants then expert second-layer codes; page-based units on paper \pc{campbell2013}{19}{the PI established a gold standard against which other coders were compared}.

\cf{Limitations and extensions} Unitization may bias reliability. The method assumes a preconceived scheme and is less suited to grounded-theory coding, though inductive codes can be deployed after development. New themes during deployment can be captured as memos (a separate grounded layer) without changing the validated code set \pc{campbell2013}{22}{Another approach is to attach a series of sepa- rate ‘‘memos’’ to various pieces of text}.

\cf{Use for this thesis} Directly relevant when the thesis has an expert first coder and a second coder with less domain knowledge. (1) \emph{Unitise first, then blind-code}: the walkthrough protocol already defines steps and screens as units, and the clean-file hand-off (units kept, codes removed) gives the second coder the same units. (2) Distinguish and report \emph{reliability} (independent) and \emph{agreement} (after negotiation), as the two give very different numbers (54\,\% vs.\ 96\,\%), and log which coder deferred. (3) Use different material in each pilot round and stop on a pre-stated criterion. (4) The proportion-agreement statistic is a weak choice for the thesis, since it ignores chance (see Hallgren, Gwet, Krippendorff). Prefer a chance-corrected coefficient, and use this paper for the negotiated-agreement procedure rather than for its statistic.
\end{papercard}
